\chapter{Considerações Finais}

Este relatório técnico apresentou a proposta estruturada para o desenvolvimento de um sistema de controle de próteses robóticas de membro superior, fundamentado na aquisição, processamento e classificação de sinais de eletromiografia de superfície (sEMG). A motivação central do projeto apoia-se na necessidade de desenvolver tecnologias assistivas funcionais, acessíveis e reprodutíveis., buscando mitigar a atual lacuna entre protótipos acadêmicos isolados e produtos comerciais de alto custo.

A arquitetura do sistema proposta demonstra viabilidade técnica ao integrar hardwares de baixo custo e alta capacidade de processamento, como os microcontroladores ESP32, a um circuito robusto de condicionamento analógico do sinal biológico. Além disso, a adoção de Redes Neurais Convolucionais (CNNs) para a extração automática de características e classificação dos gestos representa uma abordagem moderna e promissora. A estratégia metodológica de utilizar inicialmente a base de dados pública e validada NinaPro DB2 garante um ambiente seguro para o treinamento e ajuste fino dos modelos de aprendizado de máquina, reduzindo riscos de overfitting antes da etapa prática de coleta de dados com voluntários.

Por fim, o embasamento teórico consolidado e a metodologia aqui estabelecida fornecem as bases necessárias para a execução do projeto. O cronograma estipulado delineia claramente o roteiro de atividades para os próximos doze meses, englobando as disciplinas de Trabalho de Conclusão de Curso I e II. A expectativa é que, ao término deste ciclo, seja entregue um protótipo mecatrônico funcional e validado, contribuindo de forma significativa para as áreas de sistemas embarcados, inteligência artificial aplicada e engenharia de reabilitação.